\section{Experiments}
\label{sec:evaluation}

In this section, we evaluate \rec and \dic with and without parameter sharing in two multi-agent problems and compare it with a no-communication shared parameters baseline (NoComm).
Results presented are the average performance across several runs, where those without parameter sharing~(-NS), are represented by dashed lines. Across plots, rewards are normalised by the highest average reward achievable given access to the true state (Oracle).

In our experiments, we use an $\epsilon$-greedy policy with $\epsilon = 0.05$, the discount factor is $\gamma=1$, and the target network is reset every $100$ episodes. To stabilise learning, we execute parallel episodes in batches of $32$.
The parameters are optimised using RMSProp with momentum of $0.95$ and a learning rate of \num{5e-4}. %\cite{tieleman2012lecture}
The architecture makes use of \emph{rectified linear units} (ReLU), and \emph{gated recurrent units} (GRU)~\cite{cho2014properties}, which have similar performance to \emph{long short-term memory}~\cite{hochreiter1997long} (LSTM) ~\cite{chung2014empirical,jozefowicz2015empirical}. Unless stated otherwise we set $\sigma = 2$, which was found to be essential for good performance. We intent to published the source code online.
 
\subsection{Model Architecture}

\begin{wrapfigure}{r}{.4\textwidth}
    \vspace{-2.0em}
    \begin{center}
        \includegraphics[width=1\linewidth]{artwork/new/arch.pdf}
    \end{center}
    \vspace{-0.5em}
    \caption{\dic architecture.}
    \label{fig:dic_arch}
    \vspace{-1.0em}
\end{wrapfigure}

\rec and \dic share the same individual model architecture. For brevity, we describe only the \dic model here. As illustrated in Figure~\ref{fig:dic_arch}, each agent consists of a recurrent neural network (RNN), unrolled for $T$ time-steps, that maintains an internal state $h$, an input network for producing a task embedding $z$, and an output network for the $Q$-values and the messages~$m$.
%
The input for agent $a$ is defined as a tuple of $(o^a_t,m^{a'}_{t-1},u^a_{t-1}, a)$. The inputs $a$ and $u^a_{t-1}$ are passed through lookup tables, and $m^{a'}_{t-1}$ through a 1-layer MLP, both producing embeddings of size $128$. $o^a_t$ is processed through a task-specific network that produces an additional embedding of the same size. The state embedding is produced by element-wise summation of these embeddings, $z^a_t = \left( \text{TaskMLP}(o^a_t) +  \text{MLP}\lbrack |M|, 128 \rbrack (m_{t-1}) + \text{Lookup}(u^a_{t-1}) + \text{Lookup}(a) \right)$. 
We found that performance and stability improved when a batch normalisation layer~\cite{ioffe2015batch} was used to preprocess $m_{t-1}$.
$z^a_t$ is processed through a 2-layer RNN with GRUs, $h^a_{1,t} = \text{GRU}\lbrack 128, 128 \rbrack (z^a_t, h^a_{1,t-1})$, which is used to approximate the agent's action-observation history. 
Finally, the output $h^a_{2,t}$ of the top GRU layer, is passed through a 2-layer MLP $Q^a_t, m^a_t = \text{MLP}\lbrack 128, 128, (|U|+|M|) \rbrack(h^a_{2,t})$.

\subsection{Switch Riddle}
\label{subsec:switch-riddle}
    
\begin{wrapfigure}{r}{.4\textwidth}
    \vspace{-2.0em}
    \begin{center}
        \includegraphics[width=1\linewidth]{artwork/switch_vis.pdf}
    \end{center}
    \vspace{-0.5em}
    \caption{\emph{Switch:} Every day one prisoner gets sent to the interrogation room where he sees the switch and chooses from ``On'', ``Off'', ``Tell'' and ``None''. } %When the prisoner is in his cell there is nothing to do but hope for the best.}
    \label{fig:switch_vis}
    \vspace{-1.5em}
\end{wrapfigure} 
The first task is inspired by a well-known riddle described as follows:
\emph{``One hundred prisoners have been newly ushered into prison. The warden tells them that starting tomorrow, each of them will be placed in an isolated cell, unable to communicate amongst each other. Each day, the warden will choose one of the prisoners uniformly at random with replacement, and place him in a central interrogation room containing only a light bulb with a toggle switch. The prisoner will be able to observe the current state of the light bulb. If he wishes, he can toggle the light bulb. He also has the option of announcing that he believes all prisoners have visited the interrogation room at some point in time. If this announcement is true, then all prisoners are set free, but if it is false, all prisoners are executed. The warden leaves and the prisoners huddle together to discuss their fate. Can they agree on a protocol that will guarantee their freedom?''}~\cite{wu2002100}. %Unfamiliar readers are encouraged to solve the riddle before proceeding. 



\textbf{Architecture.} 
In our formalisation, at time-step $t$, agent $a$ observes ${o^a_t \in {0,1}}$, which indicates if the agent is in the interrogation room. Since the switch has two positions, it can be modelled as a 1-bit message, $m^a_t$. If agent $a$ is in the interrogation room, then its actions are $u^a_t \in \{\text{``None'',``Tell''}\}$; otherwise the only action is ``None''. The episode ends when an agent chooses ``Tell'' or when the maximum time-step, $T$, is reached. The reward $r_t$ is~$0$ unless an agent chooses ``Tell'', in which case it is $1$ if all agents have been to the interrogation room and $-1$ otherwise. Following the riddle definition, in this experiment $m^a_{t-1}$ is available only to the agent $a$ in the interrogation room. Finally, we set the time horizon $T = 4n-6$ in order to keep the experiments computationally tractable. 

\textbf{Complexity.} The switch riddle poses significant protocol learning challenges. At any time-step $t$, there are $|o|^t$ possible observation histories for a given agent, with $|o| = 3$:  the agent either is not in the interrogation room or receives one of two messages when he is. For each of these histories, an agent can chose between $4 = |U| |M| $ different options, so at time-step $t$, the single-agent policy space is $\left(|U| |M|\right)^{|o|^t}=4^{3^t}$.
The product of all policies for all time-steps defines the total policy space for an agent: $\prod 4^{3^t} = 4^{(3^{T+1}-3)/2}$, where $T$ is the final time-step.
The size of the multi-agent policy space grows exponentially in $n$, the number of agents: $4^{n{(3^{T+1}-3)/2}}$. We consider a setting where $T$ is proportional to the number of agents, so the total policy space is $4^{n 3^{O(n)}}$.
For $n = 4$, the size is $4^{88572}$.
Our approach using \dic is to model the switch as a continuous message, which is binarised during decentralised execution.


\textbf{Experimental results.} Figure~\ref{fig:exp_switch}(a) shows our results for $n=3$ agents.
All four methods learn an optimal policy in $5$k episodes, substantially outperforming the NoComm baseline. \dic with parameter sharing reaches optimal performance substantially faster than \rec. Furthermore, parameter sharing speeds both methods.
%
Figure~\ref{fig:exp_switch}(b) shows results for $n=4$ agents. \dic with parameter sharing again outperforms all other methods. In this setting, \rec without parameter sharing was unable to beat the NoComm baseline. 
These results illustrate how difficult it is for agents to learn the same protocol independently. Hence, parameter sharing can be crucial for learning to communicate. \dic-NS performs similarly to \rec, indicating that the gradient provides a richer and more robust source of information.

\begin{figure}[tb]
    \centering
    \begin{subfigure}[t]{0.32\textwidth}
        \centering
        \includegraphics[width=1\linewidth]{artwork/new/sw_3_ablation.pdf}
        \caption{Evaluation of $n=3$}
    \end{subfigure}%
    ~
    \begin{subfigure}[t]{0.32\textwidth}
        \centering
        \includegraphics[width=1\linewidth]{artwork/new/sw_4_ablation.pdf}
        \caption{Evaluation of $n=4$}
    \end{subfigure}%
    ~
    \begin{subfigure}[t]{0.32\textwidth}
        \centering
        \includegraphics[width=1\linewidth]{artwork/new/switch_stategy_3.pdf}
        \caption{Protocol of $n=3$}
    \end{subfigure}%
    \caption{\emph{Switch:} (a-b) Performance of \dic and \rec, with and without ( -NS) parameter sharing, and NoComm-baseline, for $n=3$ and $n=4$ agents.
    (c) The decision tree extracted for $n=3$ to interpret the communication protocol discovered by \dic.}
    \label{fig:exp_switch}
    \vspace{-1.0em}
\end{figure}


We also analysed the communication protocol discovered by \dic for $n=3$ by sampling $1$K episodes, for which Figure~\ref{fig:exp_switch}(c) shows a decision tree corresponding to an optimal strategy.
When a prisoner visits the interrogation room after day two, there are only two options: either one or two prisoners may have visited the room before. If three prisoners had been, the third prisoner would have  finished the game. The other options can be encoded via the ``On'' and ``Off'' position respectively.


\subsection{MNIST Games}
\label{subsec:mnist-games}

In this section, we consider two tasks based on the well known MNIST digit classification dataset~\cite{lecun1998gradient}.

\begin{wrapfigure}{r}{.55\textwidth}
    \vspace{-1.5em}
    \begin{center}
        \includegraphics[width=1\linewidth]{artwork/new/mnist_game.pdf}
    \end{center}
    \vspace{-0.5em}
    \caption{MNIST games architectures.}
    \label{fig:minst_network}
    \vspace{-1.0em}
\end{wrapfigure}

\textbf{Colour-Digit MNIST} is a two-player game in which each agent observes the pixel values of a random MNIST digit in red or green of size $2\times28\times28$, while the colour label, $c^a \in {0,1}$, and digit value, $d^a \in 0..9$, are hidden.  For each agent, reward consists of two components that are antisymmetric in the action, colour, and parity (odd, even) of the digits. Only one bit of information can be sent, so agents must agree to encode/decode either colour or parity, with parity yielding greater rewards. 
The game has two steps; in the first step, both agents send a 1-bit message, in the second step they select a binary action $u^a_2$.
The reward for each agent is $r(a)= 2 (-1)^{a^a_2 + c^a+ d^{a'}} + (-1)^{a^a_2 + d^a + c^{a'}}$ and the total cooperative reward is $r_2 = r(1) + r(2)$.present results for 5 time-steps with a 1-bit of information exchanged per step and for 2 time-steps with 4 bits exchanged per step. 

\textbf{Multi-Step MNIST} is a grayscale variant that requires agents to develop a communication protocol which integrates information across many time-steps: Each step the agents send a message, $m^a_t$, and take an action $u^a_t  \in \{0,\ldots,9\}$. Only at  the final step, $t=5$, is reward  given, $r_{5}=0.5$ for each correctly guessed digit, $u^a_{5} = d^{a'}$.
As only 1-bit is sent per step, agents must find a protocol that integrates information across the four messages they exchange (the last message is not received). The protocol can be trained using  gradients in \dic,  but also needs to have a low discretisation error.

\textbf{Architecture.}  The input processing network is a 2-layer MLP $\text{TaskMLP}\lbrack (|c| \times 28\times28), 128, 128 \rbrack(o^a_t)$. Figure~\ref{fig:minst_network} depicts the generalised setting for both games. Our experimental evaluation showed improved training time using batch normalisation after the first layer. 

\textbf{Experimental results.} Figures~\ref{fig:exp_mnist}(a) and \ref{fig:exp_mnist}(b) show that \dic substantially outperforms the other methods on both games. Furthermore, parameter sharing is crucial for reaching the optimal protocol. In multi-step MNIST, results were obtained with $\sigma = 0.5$. In this task, \rec fails to learn, while in colour-digit MNIST it fluctuates around local minima in the protocol space; the NoComm baseline is stagnant at zero.
\dic's performance can be attributed to directly optimising the messages in order to reduce the global DQN error while \rec must rely on trial and error. \dic can also optimise the message content with respect to rewards taking place many time-steps later, due to the gradient passing between agents, leading to optimal performance in multi-step MNIST. To analyse the protocol that \dic learned, we sampled $1$K episodes. Figure~\ref{fig:exp_mnist}(c) illustrates the communication bit sent at time-step $t$ by agent $1$, as a function of its input digit. Thus, each agent has learned a binary encoding and decoding of the digits. These results illustrate that differentiable communication in \dic is essential to fully exploiting the power of centralised learning and thus is an important tool for studying the learning of communication protocols. We provide further insights into the difficulties of learning a communication protocol for this task with \rec compared to \dic in the Appendix\ref{app:b}.

\begin{figure}[tb]
    \centering
    \begin{subfigure}[t]{0.36\textwidth}
        \centering
        \includegraphics[width=1\linewidth]{artwork/new/cd_many_steps.pdf}
        \caption{Evaluation of Multi-Step}
    \end{subfigure}%
    ~
    \begin{subfigure}[t]{0.36\textwidth}
        \centering
        \includegraphics[width=1\linewidth]{artwork/new/cd_extra_hard_local.pdf}
        \caption{Evaluation of Colour-Digit}
    \end{subfigure}%
    ~    
    \begin{subfigure}[t]{0.27\textwidth}
        \centering
        \includegraphics[width=1\linewidth]{artwork/new/many_steps_strategy.pdf}
        \caption{Protocol of Multi-Step}
    \end{subfigure}%    
    \caption{\emph{MNIST Games:} (a,b) Performance of \dic and \rec, with and without (-NS) parameter sharing, and NoComm, for both MNIST games.
    (c) Extracted coding scheme for multi-step MNIST.}
    \label{fig:exp_mnist}
    \vspace{-1.0em}
\end{figure}



\subsection{Effect of Channel Noise}
\begin{wrapfigure}{r}{0.373\textwidth}
    \vspace{-2.0em}
    \begin{center}
        \includegraphics[width=1\linewidth]{artwork/new/hist_comm.pdf}
    \end{center}
    \vspace{-0.5em}
    \caption{DIAL's learned activations with and without noise in \magicbox.}
    \label{fig:hist_comm_noise}
    \vspace{-1.5em}
\end{wrapfigure}

The question of why language evolved to be discrete has been studied for centuries, see e.g., the overview in~\cite{studdertkennedy05discrete}. Since \dic learns to communicate in a continuous channel, our results offer an illuminating perspective on this topic. 


In particular, Figure~\ref{fig:hist_comm_noise} shows that, in the switch riddle, \dic without noise in the communication channel learns centred activations.  By contrast, the presence of noise forces messages into two different modes during learning. Similar observations have been made in relation to adding noise when training document models  \cite{hinton2011discovering} and performing classification~\cite{courbariaux2016binarynet}. In our work, we found that adding noise was essential for successful training. More analysis on this is provided in Appendix~\ref{app:c}.

